%More details on muon reconstruction can be found in Ref.~\cite{AN-15-277}.
Muons can be reconstructed as global muon which exploits the information from muon track in muon system and a tracker track.
It can also be reconstructed by matching the tracker track with its track stubs (i.e. with at least one the muon track segment created by CSC or DT hits).

As per requirement of analysis, we define {\bf loose muons} as the muons that satisfy  
$p_T > 5$, $|\eta| < 2.4$, $dxy< 0.5$ cm, $dz < 1$ cm, where $dxy$ and $dz$ are 
defined w.r.t. the PV and using the 'muonBestTrack'. Muons have to be 
reconstructed by either the Global Muon or Tracker Muon algorithm. Standalone 
Tracks of muon reconstructed only in the muon system are rejected.
Muons with \verb|muonBestTrackType==2| (standalone) are discarded even if they 
are marked as global or tracker muons.
More details on the reconstruction of muons can be seen in Ref.~\cite{AN-15-277}.


Loose muons with $\pt$ below 200 ~\GeV are considered identified muons if they 
also pass the PF muon ID (note that the naming 
convention used for these IDs differs from the muon POG naming scheme, in which
the ``tight ID'' used here is called the ``loose ID''). Loose muons with $\pt$ 
above 200 ~\GeV ~are considered identified muons if they pass the particle flow (PF) ID or the Tracker
High-$\pt$ ID, the definition of which is shown in Table \ref{tab:highPtID}.
This relaxed definition is used to increase signal efficiency for the high-mass
search. When a very heavy resonance decays to two $\cPZ$ bosons, both bosons
will be very boosted. In the lab frame, the leptons originating from decay of
a highly boosted $Z$ will be nearly collinear, and the PF ID loses 
efficiency for muons separated by approximately $\Delta R < 0.4$, which roughly 
corresponds to muons originating from $\cPZ$ bosons with $\pt > 500\GeV$.

\begin{table}[h]
    \begin{small}
    \begin{center}
    \caption{
      The requirements for a muon to pass the Tracker High-$\pt$ ID. Note that
      these are equivalent to the Muon POG High-$\pt$ ID with the global track 
      requirements removed.
      }
    \begin{tabular}{|l|l|}
      \hline
      Plain-text description         & Technical description                 \\
      \hline
      Muon station matching          & Matching of muon to the   \\
                                     & segments in two muon stations at least \\
                                     & \textbf{NB: It means that muon is} \\
                                     & \textbf{an arbitrated tracker muon.}  \\
      \hline                                                          
      Good $\pt$ measurement         & $\frac{\pt}{\sigma_{\pt}} < 0.3$      \\
      \hline
      Vertex compatibility ($x-y$)   & $d_{xy} < 2$~mm                       \\
      \hline
      Vertex compatibility ($z$)     & $d_{z} < 5$~mm                        \\
      \hline
      Pixel hits                     & One pixel hit at least               \\
      \hline
      Tracker hits                   & Hits in at least six tracker layers   \\
      \hline
    \end{tabular}
    \label{tab:highPtID}
    \end{center}
    \end{small}
\end{table}

It may happen that a single muon is reconstruted incorrectly as two muons or even more. To cope with this situation, additional ``ghost-cleaning'' is performed.  Those tracker Muons which are not Global Muons, are needed to be arbitrated. If two muons are sharing their segments upto 50\% or more, lower quality muon is removed.

